\chapter{Polygon Partitioning}
\label{ch:polygon_partitioning}

\section{Convex Decomposition}
\label{sec:convex_decomposition}

Convex decomposition\index{convex decomposition} is the process of partitioning a non-convex (concave) polygon into a set of disjoint convex polygons. This is a fundamental operation in computational geometry, as many algorithms work more efficiently on convex shapes.

\subsection{Definitions}

\begin{definition}[Convex Polygon]\label{def:convex_polygon}
A polygon $P$ is \emph{convex} if for every pair of points $a, b \in P$, the line segment $\overline{ab}$ lies entirely within $P$.
\end{definition}

\begin{definition}[Diagonal]\label{def:diagonal}
A \emph{diagonal} of a polygon is a line segment connecting two non-adjacent vertices that lies entirely within the interior of the polygon.
\end{definition}

\subsection{Theory}

The implementation uses a \textbf{triangle merging}\index{triangle merging} approach.
\begin{enumerate}
\item Triangulate the polygon (e.g., using ear clipping~\cite{meisters1975}).
\item For each shared diagonal between two adjacent triangles/polygons, check if removing the diagonal results in a convex polygon.
\item If it does, merge the two parts.
\end{enumerate}

\begin{theorem}[Hertel--Mehlhorn Bound]\label{thm:hertel_mehlhorn}
Any convex decomposition produced by the Hertel--Mehlhorn triangle-merging heuristic has at most $4k^*$ parts, where $k^*$ is the optimal (minimum) number of convex parts.
\end{theorem}

\subsection{Pseudocode}

\begin{algorithm}[htbp]
\caption{Convex Decomposition via Triangle Merging}
\label{alg:convex_decomposition}
\begin{algorithmic}[1]
\Function{ConvexDecomposition}{$P$}
  \State $T \gets \text{Triangulate}(P)$
  \State $D \gets$ list of all interior diagonals of $T$
  \For{each $d$ in $D$}
    \State Let $P_1, P_2$ be the parts sharing diagonal $d$
    \If{$\text{Merged}(P_1, P_2)$ is convex at endpoints of $d$}
      \State Merge $P_1$ and $P_2$ into a single part
    \EndIf
  \EndFor
  \State \Return final parts
\EndFunction
\end{algorithmic}
\end{algorithm}

\subsection{Parameters}

\begin{itemize}
\item \textbf{Triangulation Algorithm}: Ear clipping is used for simplicity.
\end{itemize}

\subsection{Complexity}

\begin{itemize}
\item \textbf{Time}: $O(n^2)$ for simple merging heuristics.
\item \textbf{Space}: $O(n)$.
\end{itemize}

\subsection{Usages}

Implemented in \texttt{compgeom.polygon.polygon\_decomposer.convex\_decompose\_polygon}. Used in physics engines, motion planning, and collision detection.

\subsection{Testing and Edge Cases}

\begin{itemize}
\item \textbf{Testing}: Verify all resulting parts are convex. Check that the union of parts matches the original area.
\item \textbf{Edge Cases}: Polygons with holes, self-intersecting polygons (unsupported), highly narrow ``snaky'' polygons.
\end{itemize}

\subsection{References}

\begin{itemize}
\item Hertel, S., \& Mehlhorn, K. (1983). ``Fast triangulation of simple polygons''. FCT.
\item Wikipedia: Polygon decomposition.
\end{itemize}


\section{Convex Partitioning}
\label{sec:convex_partitioning}

Convex partitioning is the problem of dividing a non-convex (concave) polygon into the \textbf{minimum possible number} of convex sub-polygons. While a simple decomposition (like triangulation) results in $n-2$ convex parts, an optimal partition typically results in significantly fewer. This problem is solvable in polynomial time for simple polygons, making it a powerful tool for geometric simplification.

\subsection{Definitions}

\begin{definition}[Simple Polygon]\label{def:simple_polygon}
A polygon is \emph{simple} if it has no self-intersections and no holes.
\end{definition}

\begin{definition}[Reflex Vertex]\label{def:reflex_vertex}
A vertex $v$ of a polygon is \emph{reflex}\index{reflex vertex} if the interior angle at $v$ exceeds $\pi$ (i.e., $180^\circ$). Reflex vertices are the source of non-convexity.
\end{definition}

\begin{definition}[Optimal Convex Partition]\label{def:optimal_partition}
An \emph{optimal convex partition} of a polygon $P$ is a decomposition into the minimum number $k^*$ of convex sub-polygons whose union equals $P$ and whose interiors are pairwise disjoint.
\end{definition}

\subsection{Theory}

The minimum number of convex parts is determined by how reflex vertices are ``resolved'' by diagonals. Each reflex vertex must be incident to at least one diagonal that splits its reflex angle.

The \textbf{Keil--Snoeyink algorithm}\index{Keil--Snoeyink algorithm} uses \textbf{dynamic programming}\index{dynamic programming} to find the optimal partition of a simple polygon in $O(n^3)$ time (or $O(r^2 n)$ where $r$ is the number of reflex vertices).
\begin{enumerate}
\item Define $DP[i][j]$ as the minimum number of convex parts in the sub-polygon formed by vertices $V_i, \dots, V_j$ and the diagonal $V_i V_j$.
\item The algorithm systematically builds up solutions for larger sub-polygons by combining solutions for smaller ones and checking if the resulting combination is convex.
\item The optimal solution handles cases where multiple reflex vertices are resolved by a single diagonal.
\end{enumerate}

\begin{theorem}[Lower Bound on Convex Parts]\label{thm:convex_parts_lower}
Any convex partition of a simple polygon with $r$ reflex vertices requires at least $\lceil r/2 \rceil + 1$ convex parts.
\end{theorem}
\begin{proof}
Each reflex vertex must be ``cut'' by at least one diagonal to reduce its interior angle below $\pi$. A single diagonal can resolve at most two reflex vertices (one at each endpoint). Thus at least $\lceil r/2 \rceil$ diagonals are needed, producing at least $\lceil r/2 \rceil + 1$ parts.
\end{proof}

\subsection{Pseudocode}

\begin{algorithm}[htbp]
\caption{Minimum Convex Partition via Dynamic Programming}
\label{alg:convex_partition}
\begin{algorithmic}[1]
\Function{MinimumConvexPartition}{polygon}
  \State $n \gets \text{len}(polygon.vertices)$
  \State $reflex \gets [v \mid v \in polygon \wedge \text{IsReflex}(v)]$
  \If{$reflex$ is empty}
    \State \Return $[polygon]$
  \EndIf
  \State $dp \gets \text{array}[n][n]$
  \For{$length = 2$ \textbf{to} $n - 1$}
    \For{$i = 0$ \textbf{to} $n - length - 1$}
      \State $j \gets i + length$
      \If{$\neg\, \text{IsDiagonal}(i, j, polygon)$}
        \State $dp[i][j] \gets \infty$
        \State \textbf{continue}
      \EndIf
      \State $dp[i][j] \gets \min_{k \in (i+1,\ldots,j)} \bigl(dp[i][k] + dp[k][j]\bigr)$
      \If{$\text{ResolvesReflex}(i, j)$}
        \State $dp[i][j] \gets dp[i][j] - 1$
      \EndIf
    \EndFor
  \EndFor
  \State \Return $\text{ExtractPartition}(dp, polygon)$
\EndFunction
\end{algorithmic}
\end{algorithm}

\subsection{Parameters}

\begin{itemize}
\item \textbf{Polygon Type}: This optimal algorithm is for \textbf{simple polygons}. For polygons with holes, the problem is NP-hard.
\item \textbf{Goal}: This specifically minimizes the number of \textbf{parts}. Other algorithms might minimize the \textbf{total length} of diagonals (Minimum Weight Convex Partition).
\end{itemize}

\subsection{Complexity}

\begin{itemize}
\item \textbf{Time Complexity}: $O(n^3)$ for the general dynamic programming approach. More optimized versions can reach $O(r^2 n)$.
\item \textbf{Space Complexity}: $O(n^2)$ to store the DP table.
\end{itemize}

\subsection{Usages}

\begin{itemize}
\item \textbf{Path Planning}: Breaking a complex floor plan into the fewest convex ``rooms'' to simplify robot navigation.
\item \textbf{Computer Vision}: Representing a complex silhouette as a small number of convex components for shape matching.
\item \textbf{Physics Engines}: Reducing the number of primitives needed for collision detection while maintaining accuracy.
\item \textbf{Pattern Recognition}: Identifying structural features in complex polygonal datasets.
\end{itemize}

\begin{example}[Optimal Partition of an L-Shape]\label{ex:l_shape_partition}
Consider an L-shaped polygon with vertices (in order):
\[
(0,0),\; (3,0),\; (3,1),\; (1,1),\; (1,3),\; (0,3).
\]
This polygon has $r = 1$ reflex vertex at $(1,1)$. By the lower bound (\cref{thm:convex_parts_lower}), at least $\lceil 1/2\rceil + 1 = 2$ convex parts are needed. A single diagonal from $(1,1)$ to $(0,0)$ or from $(1,1)$ to $(3,3)$ would resolve the reflex vertex, yielding exactly 2 convex parts. This is optimal.
\end{example}

\subsection{Testing and Edge Cases}

\begin{itemize}
\item \textbf{Convex Polygon}: Should return the original polygon (1 part).
\item \textbf{Star Polygon}: An optimal partition might be smaller than the number of reflex vertices.
\item \textbf{``C'' Shape}: Should be split into 3 convex parts optimally.
\item \textbf{``L'' Shape}: Should be split into 2 parts.
\item \textbf{Numerical Precision}: Correctness of \texttt{IsDiagonal} and \texttt{IsReflex} is critical.
\item \textbf{Holes}: Verify the algorithm correctly identifies that holes make the optimal solution much harder.
\end{itemize}

\subsection{References}

\begin{itemize}
\item Keil, J. M., \& Snoeyink, J. S. (2002). ``On the time complexity of optimal convex polygon partitioning''. Computational Geometry.
\item Hertel, S., \& Mehlhorn, K. (1983). ``Fast triangulation of simple polygons''. FCT.
\item Chazelle, B., \& Dobkin, D. P. (1985). ``Optimal convex decompositions''. Computational Geometry.
\item Wikipedia: Convex polygon.
\end{itemize}
